\section{Limitations}
\label{sec:limitations}

\paragraph{Corpus size and homogeneity.}
At 82 documents and 3{,}134 chunks the corpus is a realistic size for an edge
deployment, but far below that of a main-conference retrieval benchmark, so the
results are best read as a case study. The documents themselves are uniform:
they belong to five business categories and share vocabulary, layout and
boilerplate. The negative result for BM25 may well be a property of this
homogeneity, and the size of the recall gains may not transfer either.
Expanding the corpus to several hundred documents is left for future work.

\paragraph{Synthetic queries.}
The query set was generated by two language models, paraphrased and judged by
one of the two, and filtered automatically; a manual spot check of twenty pairs
found seventeen clearly correct, two debatable and one mismatch. Absolute metric
values are therefore optimistic by roughly five to fifteen percent. Every
comparison in this paper is made on the same query set and is unaffected by the
bias, but the values should not be taken as estimates of performance on
naturally occurring questions, since synthetic and authentic query distributions
are known to diverge in ways that aggregate scores do not reveal
\citep{kucia2026synthetic}. Because the judge is also one of the generators, the
acceptance filter shares the generator's blind spots.

\paragraph{Statistical power.}
The reranked arms are measured on sixty queries, which bounds the resolution of
the quality comparison: the bootstrap interval for the gain of the reranker
over the contextualized dense-only arm includes zero
(Section~\ref{sec:results-main}), so this study cannot decide whether a 0.6B
reranker adds quality here. The claim we defend is narrower. The reranker is not
required for the pipeline to work, because the distilled reranker-free
configuration is indistinguishable from the reranked one at three orders of
magnitude lower latency, and because reranking flat chunks does not recover the
missing context. Larger query sets would be needed to resolve what remains.

\paragraph{Ablation methodology.}
The dimension and instruction ablations reuse the indexed 1{,}024-dimensional
vectors and truncate them after the fact. Truncating a Matryoshka embedding is
equivalent to evaluating the corresponding prefix, but a model trained natively
at 512 dimensions could behave differently. The fusion ablations likewise reuse
the dense candidate lists and vary only the fusion parameters.

\paragraph{Reranker configuration.}
The reranked arms use the harness default of sixty-four candidates. An ablation
over pools of 16, 32 and 64 on an earlier sixty-query subset shows identical
recall@10, a small loss in recall@5 at 64, where the reranker begins to reorder
correct pairs, and a latency that is not monotone (42.9~s, 36.6~s and 44.0~s at
pools of 16, 32 and 64). Latency here is dominated by per-call overhead more
than by pool size, so a smaller pool is the better trade at query time, and the
reported arm~\arm{c} and arm~\arm{d} numbers are conservative. We did not
evaluate alternative rerankers; a stronger cross-encoder might change the
quality comparison, though not the latency ratio.

\paragraph{Language and domain.}
The corpus is Russian and the tasks are document-centric. That a Russian query
instruction is at least as good as an English one is specific to this setting,
and it is not a general rule for multilingual retrieval.

\paragraph{No end-to-end evaluation.}
We measure retrieval quality, not answer quality. Word-level F1 against the
teacher context measures how faithfully the student reproduces the teacher, and
the retrieval metrics measure whether the correct chunk is found; neither
measures whether a downstream generator produces a correct answer. A downstream
evaluation could reveal that shorter student contexts are less useful than their
retrieval metrics suggest.

\paragraph{Hardware.}
All measurements come from a single Apple~M2 with MLX, the retrieval models in
fp16 and the ingest model in bfloat16. Absolute latency and throughput will
differ on other devices. The ratio between a 0.6B cross-encoder over a candidate
pool and a single dense search is large enough to hold on any hardware that can
run both, but the exact factor will not.
